% Source: https://artofproblemsolving.com/community/c3h1717416_junior_regional_olympiad__fbh_2014_grade_7_problem_4
% Stable identity: aops-topic-1717416-q01
The remainders satisfy $0\le a\le2$, $0\le b\le4$, $0\le c\le6$, so $4a+3b+2c\le8+12+12=32$.
\begin{itemize}
\item $a=2$: $3b+2c=22$. Since $2c\le12$, $3b\ge10$, so $b=4$ and $c=5$.
\item $a=1$: $3b+2c=26$ needs $3b\ge14$, so $b\ge5$: impossible.
\item $a=0$: $3b+2c=30$ needs $3b\ge18$: impossible.
\end{itemize}
So $n\equiv2\pmod3$, $n\equiv4\pmod5$, $n\equiv5\pmod7$. The first two say $n\equiv-1\pmod{15}$, so $n$ is one of $14,29,44,59,74,89,104$ modulo $105$; their remainders modulo $7$ are $0,1,2,3,4,5,6$, and only $89$ leaves remainder $5$. Hence the remainder is $\boxed{89}$.
